%% kapa: exam=1 round=27 year=2016 period=1 form=A subject=economics subject_ko=경제학원론 seq=25 number=65
%% source: Q-Net 감정평가사 기출문제 (cst00309) artlSeq=5169402 file=extracted/1008/5169402/제27회 감정평가사 1차 문제/제27회 감정평가사 1교시문제(A형).pdf
%% conversion: kordoc 4.24.51 -> tools/gen1.py (auto); review=proofread-sonnet; flags=none
\begin{question}{65}
A국의 거시경제모형이 다음과 같을 때, 총수요곡선으로 옳은 것은?
\begin{framedtext}
민간소비: $C = 2 + 0.5Y$ \par
투자: $I = 2 - r$ \par
정부지출: $G = 3$ \par
실질화폐수요: $\dfrac{M^D}{P} = 4 + 0.5Y - r$ \par
명목화폐공급: $M^S = 3$ \par
(단, $r$은 이자율, $Y$는 국민소득, $P$는 물가수준이고, $Y>3$이다.) \par
\end{framedtext}
\begin{choices}
  \item $Y = 1 + \dfrac{1}{P}$
  \item $Y = 2 + \dfrac{2}{P}$
  \item $Y = 3 + \dfrac{3}{P}$
  \item $Y = 4 + \dfrac{4}{P}$
  \item $Y = 5 + \dfrac{5}{P}$
\end{choices}
\end{question}
\answer{3}
